\newpage
\section*{Appendix B\quad Reproducing every number}
\addcontentsline{toc}{section}{B Reproducing every number}
\newcommand{\nElevenScriptCount}{42}\newcommand{\nElevenAssertCount}{98}


These notes make one promise: no number in them was typed by hand. How would you check a promise like
that? This appendix tells you how. Every number comes from a script in the \texttt{code/} folder of the companion
repository; every script writes its numbers into a small file of \LaTeX\ macros, and the text only ever
uses those macros. If you change a script and rebuild, the text changes with it.

\subsection*{B.1\quad Rebuilding everything}

From the top folder of the repository:
\begin{center}
\small
\begin{tabular}{@{}P{6.2cm} P{7.6cm}@{}}
\toprule
Command & What it does\\
\midrule
\texttt{make} & runs every number script whose source changed, then builds the PDF with \texttt{latexmk}\\
\texttt{make clean} & deletes the generated macro files and the build, so the next \texttt{make} redoes everything\\
\texttt{python3 lint/lint.py tex/ch*.tex} & runs the style checks of the design pack on every chapter\\
\texttt{python3 code/ch11\_ablate.py --retrain} & retrains the Chapter 11 models (about 5 CPU minutes)\\
\texttt{python3 code/ch11\_appB\_index.py} & regenerates the script table below from the code itself\\
\bottomrule
\end{tabular}
\end{center}
The slow step, training the small models of Chapters 8 and 11, caches its results in \texttt{tex/data\_ch8/}
and \texttt{tex/data\_ch11/}, so an ordinary rebuild takes about a minute. To check that a rebuild is
faithful, run \texttt{make clean}, retrain, rebuild, and compare the two PDFs' text: with the fixed seeds
the numbers should agree to the digits printed, except the training times.

\subsection*{B.2\quad Every script, what it writes, and what it checks}

The table is generated from the scripts themselves. ``Checks'' counts the \texttt{assert} statements
in each script: places where the script stops the build if something disagrees, for example a
hand-calculation against NumPy, or our parameter calculator
against the real Hugging Face model classes. In total there are \nElevenScriptCount\ scripts and
\nElevenAssertCount\ checks.

\begin{center}
\scriptsize
\begin{longtable}{@{}P{3.0cm} P{7.2cm} P{3.2cm} r@{}}
\toprule
Script & Purpose & Writes & Checks\\
\midrule
\endhead
\csname @@input\endcsname data_ch11/scripts.tex
\bottomrule
\end{longtable}
\end{center}

\paragraph{Checks you can rely on.} Ask yourself, for any number you doubt: what would have caught it if
it were wrong? A few checks deserve a sentence each.
\begin{itemize}
\item The toy attention of Chapter 3 is computed twice, once with NumPy and once with exact fractions,
and the two must agree to $10^{-12}$.
\item The parameter calculator of Chapters 10 and 11 must equal the parameter count of the real model
classes for GPT-2 (\nElevenGptParams) and Qwen3-0.6B (\nElevenQwenParams), built without downloading any
weights; two calculators written separately must agree.
\item Derivatives in Appendix A are checked against small finite differences.
\item Every chapter must pass \texttt{lint/lint.py} with no hits.
\end{itemize}

\subsection*{B.3\quad Software}

\begin{center}
\small
\begin{tabular}{@{}l l@{}}
\toprule
Tool & Version used for this edition\\
\midrule
\csname @@input\endcsname data_ch11/versions.tex
\bottomrule
\end{tabular}
\end{center}
Training runs used one CPU thread (\texttt{torch.set\_num\_threads(1)}) and fixed seeds. A different
machine can give slightly different losses in the last digit; the conclusions in \S11.6 are stated
against a noise level measured with two seeds for exactly this reason.

\subsection*{B.4\quad What needs the network, and what could not be done here}

Building the notes needs no network at all. Three things would, and in the environment where this
edition was made they were blocked:
\begin{itemize}
\item \textbf{Hugging Face} (\texttt{huggingface.co}): the published \texttt{config.json}, tokenizer and
weights of Qwen3-0.6B. Until it is reachable, the Qwen3 configuration comes from the Qwen3 technical report
(2025) and is marked as such in Chapter 10. Our parameter count for that configuration matches the real
model class exactly, which checks our arithmetic but not the configuration itself.
\item \textbf{Wikipedia} (\texttt{en.wikipedia.org}, \texttt{dumps.wikimedia.org}): the live page
\emph{Signs of AI writing}, against which the lint rules must be re-checked, and Simple English Wikipedia
text as a second corpus for Chapters 1 and 2. Until then, the corpus numbers come from the talk corpus
only: the nine parts of a corpus of transcribed technical talks on AI engineering
(\texttt{data/ai\_week1.txt} to \texttt{data/ai\_week9.txt}), parts 1--8 for training and part 9 held out.
\item \textbf{Large models}: anything that needs a GPU (running Qwen3-0.6B itself, longer training) was
out of scope. The scripts are written so that the same code runs with a larger budget.
\end{itemize}

\subsection*{B.5\quad How we will know the notes teach}

Numbers that regenerate show that the notes are \emph{consistent}. Do they show that the notes
\emph{teach}? No. For that we need readers, and the plan before publication is:
\begin{enumerate}
\item \textbf{Pre-test and post-test} with 5 to 10 readers, at least two of whom have been out of
education for a long time. Each test uses questions written to the chapter objectives. We report the
normalised gain (Hake, 1998): the share of the possible improvement that each reader actually made.
\item \textbf{Item analysis} of every \emph{Check yourself} question: how many readers got it right
(questions everyone gets right or wrong teach us nothing), whether it separates stronger from weaker
readers, and whether every wrong option was chosen by at least one reader. A wrong option nobody chooses
means the misconception behind it was not real, and the option is replaced.
\item \textbf{Think-aloud readings} with 3 to 5 readers who say out loud what they are thinking as they
read (Ericsson and Simon, 1993). Every place where a reader stalls, rereads or guesses wrongly goes into
a list, and every item on the list is fixed and tested again.
\end{enumerate}
Which of these would convince you, if you were the reader? The results, and the changes they cause, will be published with the next edition alongside an errata
page.
